% Options for packages loaded elsewhere
\PassOptionsToPackage{unicode}{hyperref}
\PassOptionsToPackage{hyphens}{url}
\documentclass[
]{article}
\usepackage{xcolor}
\usepackage{amsmath,amssymb}
\setcounter{secnumdepth}{-\maxdimen} % remove section numbering
\usepackage{iftex}
\ifPDFTeX
  \usepackage[T1]{fontenc}
  \usepackage[utf8]{inputenc}
  \usepackage{textcomp} % provide euro and other symbols
\else % if luatex or xetex
  \usepackage{unicode-math} % this also loads fontspec
  \defaultfontfeatures{Scale=MatchLowercase}
  \defaultfontfeatures[\rmfamily]{Ligatures=TeX,Scale=1}
\fi
\usepackage{lmodern}
\ifPDFTeX\else
  % xetex/luatex font selection
\fi
% Use upquote if available, for straight quotes in verbatim environments
\IfFileExists{upquote.sty}{\usepackage{upquote}}{}
\IfFileExists{microtype.sty}{% use microtype if available
  \usepackage[]{microtype}
  \UseMicrotypeSet[protrusion]{basicmath} % disable protrusion for tt fonts
}{}
\makeatletter
\@ifundefined{KOMAClassName}{% if non-KOMA class
  \IfFileExists{parskip.sty}{%
    \usepackage{parskip}
  }{% else
    \setlength{\parindent}{0pt}
    \setlength{\parskip}{6pt plus 2pt minus 1pt}}
}{% if KOMA class
  \KOMAoptions{parskip=half}}
\makeatother
\usepackage{longtable,booktabs,array}
\newcounter{none} % for unnumbered tables
\usepackage{calc} % for calculating minipage widths
% Correct order of tables after \paragraph or \subparagraph
\usepackage{etoolbox}
\makeatletter
\patchcmd\longtable{\par}{\if@noskipsec\mbox{}\fi\par}{}{}
\makeatother
% Allow footnotes in longtable head/foot
\IfFileExists{footnotehyper.sty}{\usepackage{footnotehyper}}{\usepackage{footnote}}
\makesavenoteenv{longtable}
\setlength{\emergencystretch}{3em} % prevent overfull lines
\providecommand{\tightlist}{%
  \setlength{\itemsep}{0pt}\setlength{\parskip}{0pt}}
\usepackage{bookmark}
\IfFileExists{xurl.sty}{\usepackage{xurl}}{} % add URL line breaks if available
\urlstyle{same}
% fallback for those not using the hyperref driver hyperxmp:
\makeatletter
\@ifundefined{xmpquote}{\newcommand{\xmpquote}[1]{#1}}{}
\makeatother
\hypersetup{
  hidelinks,
  pdfcreator={LaTeX via pandoc}}

\author{}
\date{}

\begin{document}

\section{Measuring the Defenders: A Layer-Aware, Framework-Mapped
Benchmark for Model Context Protocol Security
Proxies}\label{measuring-the-defenders-a-layer-aware-framework-mapped-benchmark-for-model-context-protocol-security-proxies}

\textbf{Author:} Gowthaman Arumugam, Independent Researcher
(agowthaman90@gmail.com) \textbf{Version:} 0.3 · 2026-07-14 ·
\emph{preprint, not peer-reviewed} \textbf{Artifacts:} mcp-defense-bench
(benchmark) · mcp-bastion (reference proxy) --- see
\hyperref[availability]{Availability} \textbf{Changes in v0.3:} two
newly-published (June 2026) attacks added as vectors --- mid-session
tool injection (MSTI) and multi-tool split poisoning (ShareLock) ---
bringing the rubric to 24 vectors / 33 cases; the crosswalk gains a
third U.S.-government framework column, the NSA MCP Security guidance
(May 2026); mcp-bastion adds cross-tool correlation and now detects
ShareLock (coverage 34\% → 38\%); and we add a tool-class taxonomy
separating content-scanning proxies (scored) from
access-control/isolation gateways (tracked).

\begin{center}\rule{0.5\linewidth}{0.5pt}\end{center}

\subsection{Abstract}\label{abstract}

The Model Context Protocol (MCP) has become a common way to connect
large-language-model (LLM) agents to external tools and data, and a
substantial 2025--2026 literature now documents its attack surface.
Existing security benchmarks for MCP measure how well an \textbf{LLM
agent resists} attacks. None, to our knowledge, measure how well a
\textbf{defensive proxy or gateway covers} that attack surface, and none
map their results to the governance frameworks organizations actually
use --- the NIST AI Risk Management Framework (AI RMF) and the OWASP Top
10 for LLM and Agentic Applications. We present two artifacts to fill
this gap: (1) a \textbf{threat--control crosswalk} of 24 MCP attack
vectors mapped to architectural layer, STRIDE class, NIST AI RMF
function, OWASP category, and the NSA MCP Security guidance; and (2) an
open, vendor-neutral \textbf{defense-side benchmark} that runs a
candidate MCP security proxy against a corpus of attack fixtures with
matched benign controls, and scores its coverage using the crosswalk as
a rubric. Each tool is driven through its \emph{real} code, not a mock.
We evaluate three open-source proxies (mcp-bastion, mcp-firewall,
pipelock) plus a null baseline over 33 test cases. The strongest proxy
reaches \textbf{38\% weighted coverage (15 of 24 vectors)} --- a proxy's
realistic ceiling --- and even the union of all three tools leaves
\textbf{9 of 24 vectors covered by none}: the entire
registry/supply-chain surface, OS-isolation-dependent vectors, and
several semantic vectors. All results carry zero false positives against
matched benign controls. We argue this is direct, quantified evidence
that a runtime proxy alone is insufficient and that defense-in-depth
spanning \emph{distinct mechanism classes} (proxy, cryptographic
attestation, OS isolation) is required. We also report two
methodological findings we observed first-hand: benchmark scores are
sensitive to how attack fixtures are encoded (a fairness hazard), and
the benchmark can \emph{guide} a defender's development --- over the
study, mcp-bastion's coverage rose from 9\% to 38\% as gaps it surfaced
were fixed and re-verified --- including two attacks first published in
June 2026, one of which (ShareLock) no measured tool covered until a
cross-tool correlation check was added. An evasion-robustness dimension
makes the encoding-sensitivity point concrete: obfuscated variants
(zero-width, homoglyph, base64) of known attacks are caught by
\emph{different} tools, and no single tool survives all three. A matched
benign-control design and a ``measure only what the tool inspects at
runtime'' integrity rule keep the results honest.

\begin{center}\rule{0.5\linewidth}{0.5pt}\end{center}

\subsection{1. Introduction}\label{introduction}

MCP standardizes how an LLM agent discovers and calls external ``tools''
exposed by ``servers.'' Its rapid adoption across consumer and
enterprise agents has been accompanied, in 2025--2026, by a dense body
of security research: systematization-of-knowledge papers, formal threat
taxonomies, attack benchmarks, and cryptographic attestation proposals
(Section 2). This work establishes \emph{what can go wrong}.

A practical question remains under-addressed: \textbf{given a defensive
tool that sits in front of MCP servers --- a proxy, gateway, or scanner
--- how much of the known attack surface does it actually cover?}
Answering it well requires three things that, individually, exist in the
literature but have not been combined:

\begin{enumerate}
\def\labelenumi{\arabic{enumi}.}
\tightlist
\item
  a \textbf{shared taxonomy} of MCP attack vectors;
\item
  a mapping from those vectors to the \textbf{governance frameworks}
  (NIST AI RMF, OWASP) that security and procurement teams reason in;
  and
\item
  an \textbf{executable, reproducible} way to test a defender and report
  coverage.
\end{enumerate}

We contribute all three, plus an empirical study. Concretely:

\begin{itemize}
\tightlist
\item
  \textbf{A threat--control crosswalk} (Section 3): 24 MCP attack
  vectors, each mapped to one or more of six architectural layers, a
  STRIDE class, NIST AI RMF functions, NSA MCP Security recommendation
  areas, and OWASP LLM-2025 and Agentic-2026 categories. Published as
  open, versioned, machine-readable data.
\item
  \textbf{A defense-side benchmark} (Section 4): for each vector, one or
  more attack \textbf{fixtures} paired with a near-identical
  \textbf{benign control}; a small \textbf{adapter} per tool that drives
  the tool's real detection code; and a scorer that reports coverage per
  layer and per framework.
\item
  \textbf{An empirical evaluation} (Sections 5--6) of three open-source
  proxies and a null baseline, yielding the complementarity finding
  above.
\item
  \textbf{A candid methodology discussion} (Sections 7--8): the
  corpus-encoding sensitivity we observed, the matched-control and
  runtime-integrity safeguards, and the limits of the study.
\end{itemize}

Our aim is not to rank products but to map coverage, expose gaps, and
provide reusable infrastructure.

\subsection{2. Related Work}\label{related-work}

\begin{quote}
All 2603/2604-series works below are recent, non-peer-reviewed preprints
whose taxonomies are their authors' own constructs; we cite them as
\emph{proposed by}, not as established consensus. Coverage statistics
quoted from any single paper reflect that paper's surveyed set.
\end{quote}

\textbf{Threat taxonomies.} Recent work systematizes MCP threats through
several lenses: an SoK separating adversarial security threats (indirect
prompt injection, tool poisoning) from epistemic safety hazards across
the Resources/Prompts/Tools primitives {[}SoK-2512.08290{]}; a formal
framework of 7 threat categories and 23 attack vectors grounded in a
large tool corpus {[}Formal-2604.05969{]}; a STRIDE/DREAD model over
five MCP components that identifies tool poisoning as the most impactful
client-side vulnerability {[}STRIDE-2603.22489{]}; and a
defense-placement taxonomy classifying defenses by the layer that
enforces them, which finds protection ``uneven and predominantly
tool-centric'' with weak transport, host-orchestration, and supply-chain
coverage {[}MCP-DPT-2604.07551{]}. Our crosswalk reuses these vector
families but adds the framework mapping none of them provides.

\textbf{Attack-side benchmarks.} MSB {[}MSB-2510.15994{]} (12 attacks,
\textasciitilde2,000 instances, real MCP execution) measures the
resistance of the \textbf{LLM agent} across the tool-use pipeline;
MCPTox {[}MCPTox-2508.14925{]} scopes to tool poisoning;
AgentDefense-Bench aggregates datasets into MCP-JSON-RPC test cases;
MCP-Bench {[}MCPBench-2508.20453{]} measures agent \emph{capability},
not security. All measure the agent (or the model), not a
\textbf{defensive proxy}. Ours measures the defender, on the full mapped
surface, and reports coverage crosswalked to NIST/OWASP --- which we
could not find in any surveyed benchmark.

\textbf{Attestation and provenance.} ETDI {[}ETDI-2506.01333{]} proposes
signed, versioned tool definitions; a ``Trustworthy MCP Registry''
blueprint {[}TrustReg-FI2026{]} composes existing standards (well-known
URIs, Sigstore, JSON Canonicalization + JWS) for provenance and runtime
integrity, and notes the official MCP registry remains an unverified
pointer architecture. These are complementary to our benchmark; a future
attestation-conformance suite (Section 9) would test them.

\textbf{Scanners and gateways.} Practitioner tools include Invariant
Labs' MCP-Scan (since acquired by Snyk and now a cloud-gated scanner),
Lasso's and EnkryptAI's MCP gateways, Docker's isolation-based MCP
gateway, and the proxies we evaluate here. Our benchmark scores
locally-deterministic defenders and treats cloud-model defenders as
observable-but-not-reproducibly-scoreable (Section 4.4).

\subsection{3. The Threat--Control
Crosswalk}\label{the-threatcontrol-crosswalk}

The crosswalk is the benchmark's rubric and a standalone artifact. It
enumerates \textbf{24 MCP attack vectors} consolidated from the
taxonomies in Section 2, and maps each to:

\begin{itemize}
\tightlist
\item
  \textbf{Architectural layer} (one of: host-orchestration, client,
  transport, server, tool, registry-supply-chain) --- following the
  defense-placement view of {[}MCP-DPT-2604.07551{]};
\item
  \textbf{STRIDE} class;
\item
  \textbf{NIST AI RMF} function(s): Govern, Map, Measure, Manage;
\item
  \textbf{NSA MCP Security} recommendation area (Authentication \&
  Access Control, Monitoring \& Logging, Data Classification \&
  Segmentation, Runtime Controls) --- from the NSA AI Security Center
  guidance {[}NSA-MCP-2026{]};
\item
  \textbf{OWASP Top 10 for LLM Applications (2025)} category
  (LLM01--LLM10); and
\item
  \textbf{OWASP Top 10 for Agentic Applications (2026)} category
  (ASI01--ASI10).
\end{itemize}

Table 1 lists the vectors and their primary layer. The full mapping,
including STRIDE and all framework codes, is published as
machine-readable data (\texttt{rubric/crosswalk.json}, CC-BY-4.0).

\textbf{Table 1. The 24 vectors (primary layer).}

{\def\LTcaptype{none} % do not increment counter
\begin{longtable}[]{@{}
  >{\raggedleft\arraybackslash}p{(\linewidth - 10\tabcolsep) * \real{0.1667}}
  >{\raggedright\arraybackslash}p{(\linewidth - 10\tabcolsep) * \real{0.1667}}
  >{\raggedright\arraybackslash}p{(\linewidth - 10\tabcolsep) * \real{0.1667}}
  >{\raggedleft\arraybackslash}p{(\linewidth - 10\tabcolsep) * \real{0.1667}}
  >{\raggedright\arraybackslash}p{(\linewidth - 10\tabcolsep) * \real{0.1667}}
  >{\raggedright\arraybackslash}p{(\linewidth - 10\tabcolsep) * \real{0.1667}}@{}}
\toprule\noalign{}
\begin{minipage}[b]{\linewidth}\raggedleft
\#
\end{minipage} & \begin{minipage}[b]{\linewidth}\raggedright
Vector
\end{minipage} & \begin{minipage}[b]{\linewidth}\raggedright
Layer
\end{minipage} & \begin{minipage}[b]{\linewidth}\raggedleft
\#
\end{minipage} & \begin{minipage}[b]{\linewidth}\raggedright
Vector
\end{minipage} & \begin{minipage}[b]{\linewidth}\raggedright
Layer
\end{minipage} \\
\midrule\noalign{}
\endhead
\bottomrule\noalign{}
\endlastfoot
1 & Tool poisoning & tool & 13 & Sandbox escape & server \\
2 & Tool shadowing / name collision & client & 14 & Schema / validation
bypass & server \\
3 & Rug pull (definition mutation) & tool & 15 & Man-in-the-middle
(transport) & transport \\
4 & Out-of-scope parameter injection & tool & 16 & DNS rebinding (local
servers) & transport \\
5 & Prompt injection via tool results & tool & 17 & Server impersonation
& registry-supply-chain \\
6 & Indirect / retrieval injection & tool & 18 & Excessive permission /
escalation & host-orchestration \\
7 & Cross-tool exfiltration (confused deputy) & client & 19 & Credential
/ token theft & host-orchestration \\
8 & Tool-transfer / cross-server chaining & host-orchestration & 20 &
Consent fatigue / over-broad grants & client \\
9 & False-error escalation & tool & 21 & Command injection & server \\
10 & Package / name squatting & registry-supply-chain & 22 &
System-prompt / context leakage & client \\
11 & Supply-chain poisoning (provenance gap) & registry-supply-chain &
23 & Mid-session tool injection (MSTI) & client \\
12 & Configuration drift & server & 24 & Multi-tool split poisoning
(ShareLock) & tool \\
\end{longtable}
}

The mappings are indicative and reviewable, not a certification --- we
invite community correction, which is itself a goal of publishing the
rubric openly. The OWASP Agentic categories used are the official 2026
set: ASI01 Agent Goal Hijack, ASI02 Tool Misuse and Exploitation, ASI03
Identity and Privilege Abuse, ASI04 Agentic Supply Chain
Vulnerabilities, ASI05 Unexpected Code Execution (RCE), ASI06 Memory \&
Context Poisoning, ASI07 Insecure Inter-Agent Communication, ASI08
Cascading Failures, ASI09 Human-Agent Trust Exploitation, ASI10 Rogue
Agents.

\subsection{4. Benchmark Methodology}\label{benchmark-methodology}

\subsubsection{4.1 Design}\label{design}

For each vector the corpus contains one or more \textbf{test cases}. A
test case is a static JSON object with: a \textbf{malicious fixture}
(the attack, as data --- a tool definition, tool call, tool result,
registry entry, or scenario), a near-identical \textbf{benign control},
the \textbf{expected} defender behavior
(\texttt{detect}/\texttt{enforce}), and a \textbf{provenance} note.
Fixtures are static data, never live exploits; attacker hosts use
reserved names (\texttt{attacker.example}; the cloud-metadata address
\texttt{169.254.169.254} appears only as an SSRF \emph{target} in data,
never contacted).

\subsubsection{4.2 Adapters and the runtime-integrity
rule}\label{adapters-and-the-runtime-integrity-rule}

Each evaluated tool provides a thin \textbf{adapter} exposing
\texttt{assess(input)\ →\ \{detect,\ enforce\}}. The runner calls it
once on the malicious fixture and once on the benign control. Adapters
drive the tool's \textbf{real} code (e.g., importing the proxy's
detection functions, or invoking its CLI/SDK with a committed
configuration).

A central rule keeps scores honest:

\begin{quote}
\textbf{Report only what the tool actually inspects at runtime --- never
what its rules \emph{could} match if pointed at data the tool never
sees.}
\end{quote}

For example, a proxy that inspects tool \emph{definitions} but not tool
\emph{results} returns ``not detected'' for result-borne attacks, even
if its text rules would match the result string in isolation. This
measures deployed behavior, not theoretical capability.

\subsubsection{4.3 Scoring}\label{scoring}

Per test case: if the malicious fixture is not flagged, the vector
scores \textbf{none}; if it is flagged but the benign control is
\emph{also} flagged, the vector scores \textbf{none} and is marked a
\textbf{false positive}; if the malicious fixture is flagged and the
benign control is clean, the vector scores \textbf{detect} (warn) or
\textbf{enforce} (block). A vector's level is the best across its test
cases. Weighted coverage assigns enforce = 1.0, detect = 0.5, none = 0,
over 24 vectors. The matched benign control is essential: without it, a
tool that flags everything would score perfectly while being useless.

\subsubsection{4.4 Reproducibility rule}\label{reproducibility-rule}

A defender is \textbf{scored} only if its detection logic runs
\textbf{locally and deterministically}. Tools that delegate detection to
a proprietary cloud model can be \emph{observed} but not reproducibly
\emph{scored}, since their verdicts depend on a black box that can
change between runs. This excludes, e.g., the cloud-gated Snyk
agent-scan; it includes the three proxies below.

\subsection{5. Experimental Setup}\label{experimental-setup}

We evaluate four adapters over a 33-case corpus (24 base cases + 6
realistic ``v2'' cases for fairness + 3 ``v3'' evasion-robustness cases;
see Section 7):

{\def\LTcaptype{none} % do not increment counter
\begin{longtable}[]{@{}
  >{\raggedright\arraybackslash}p{(\linewidth - 6\tabcolsep) * \real{0.2500}}
  >{\raggedright\arraybackslash}p{(\linewidth - 6\tabcolsep) * \real{0.2500}}
  >{\raggedright\arraybackslash}p{(\linewidth - 6\tabcolsep) * \real{0.2500}}
  >{\raggedright\arraybackslash}p{(\linewidth - 6\tabcolsep) * \real{0.2500}}@{}}
\toprule\noalign{}
\begin{minipage}[b]{\linewidth}\raggedright
Tool
\end{minipage} & \begin{minipage}[b]{\linewidth}\raggedright
Class
\end{minipage} & \begin{minipage}[b]{\linewidth}\raggedright
License
\end{minipage} & \begin{minipage}[b]{\linewidth}\raggedright
How driven
\end{minipage} \\
\midrule\noalign{}
\endhead
\bottomrule\noalign{}
\endlastfoot
\textbf{null-baseline} & control & --- & returns ``not detected'' for
all inputs \\
\textbf{mcp-bastion} & runtime proxy & Apache-2.0 & imports real
detection: \texttt{scanTool}, \texttt{scanText}, \texttt{scanToolSet}
(cross-tool correlation), \texttt{validateArguments},
\texttt{checkRequestedScopes}, \texttt{checkTransportSecurity},
\texttt{checkRequestOrigin}, \texttt{hashToolDefinition} \\
\textbf{mcp-firewall} & runtime proxy (Python) & AGPL-3.0 & real SDK
(\texttt{Gateway.check} / \texttt{scan\_response}) with a committed
config \\
\textbf{pipelock} & egress firewall (Go) & Apache-2.0 core & real
scanners via \texttt{explain\ -\/-json} (URL/egress) and
\texttt{mcp\ scan} (MCP-response injection), both offline and
deterministic \\
\end{longtable}
}

For AGPL/other licensed tools we \emph{run} the tool to benchmark it and
report scores; we do not redistribute their code. Each tool is pinned to
a released version and configured with its own recommended/starter
policy, committed to the repository for reproducibility.

\subsection{6. Results}\label{results}

\textbf{Table 2. Verified coverage (33 cases; weighted over 24
vectors).}

{\def\LTcaptype{none} % do not increment counter
\begin{longtable}[]{@{}
  >{\raggedright\arraybackslash}p{(\linewidth - 12\tabcolsep) * \real{0.1429}}
  >{\raggedright\arraybackslash}p{(\linewidth - 12\tabcolsep) * \real{0.1429}}
  >{\raggedleft\arraybackslash}p{(\linewidth - 12\tabcolsep) * \real{0.1429}}
  >{\raggedleft\arraybackslash}p{(\linewidth - 12\tabcolsep) * \real{0.1429}}
  >{\raggedleft\arraybackslash}p{(\linewidth - 12\tabcolsep) * \real{0.1429}}
  >{\raggedleft\arraybackslash}p{(\linewidth - 12\tabcolsep) * \real{0.1429}}
  >{\raggedleft\arraybackslash}p{(\linewidth - 12\tabcolsep) * \real{0.1429}}@{}}
\toprule\noalign{}
\begin{minipage}[b]{\linewidth}\raggedright
Tool
\end{minipage} & \begin{minipage}[b]{\linewidth}\raggedright
Class
\end{minipage} & \begin{minipage}[b]{\linewidth}\raggedleft
Weighted coverage
\end{minipage} & \begin{minipage}[b]{\linewidth}\raggedleft
enforce
\end{minipage} & \begin{minipage}[b]{\linewidth}\raggedleft
detect
\end{minipage} & \begin{minipage}[b]{\linewidth}\raggedleft
none
\end{minipage} & \begin{minipage}[b]{\linewidth}\raggedleft
False positives
\end{minipage} \\
\midrule\noalign{}
\endhead
\bottomrule\noalign{}
\endlastfoot
mcp-bastion & runtime proxy & \textbf{38\% (9.0/24)} & 3 & 12 & 9 & 0 /
33 \\
mcp-firewall & runtime proxy & \textbf{13\% (3.0/24)} & 3 & 0 & 21 & 0 /
33 \\
pipelock & egress firewall & \textbf{10\% (2.5/24)} & 1 & 3 & 20 & 0 /
33 \\
null-baseline & control & 0\% & 0 & 0 & 24 & 0 / 33 \\
\end{longtable}
}

\textbf{A proxy's ceiling is partial.} mcp-bastion, the broadest tool,
spans the definition, response, argument, transport, and
egress/least-privilege layers (poisoning, shadowing, rug-pull; response
and retrieval injection and result-borne leakage; schema/parameter
validation; plaintext-transport and DNS-rebinding defense;
sensitive-argument and over-broad-scope detection). mcp-firewall and
pipelock add \emph{enforcement} (deny/redact/SSRF-block) on the call and
egress layers where bastion only warns; pipelock additionally detects
response-borne prompt injection via its MCP-response scanner.

\textbf{Two newly-published attacks (June 2026).} We added mid-session
tool injection (MSTI, {[}WebMCP-MSTI-2606.06387{]}) and multi-tool split
poisoning (ShareLock, {[}ShareLock-2606.27027{]}) as vectors 23--24.
MSTI --- a tool that mutates its behavior after first approval,
mid-session --- is caught by mcp-bastion's definition pinning, which
re-hashes the tool on each use and \emph{enforces} on a post-pin change;
neither other tool covers it. ShareLock --- a payload split across
several tools so no single description looks malicious --- was covered
by \textbf{no} measured tool until we added \textbf{cross-tool
correlation} to mcp-bastion (Section 7): because a proxy sees a server's
whole tool set, it scans the combined descriptions and flags coordinated
\texttt{share}/\texttt{checksum}-style staging metadata across tools. It
now \emph{detects} ShareLock; the other two tools, which scan one tool
at a time, still miss it. This is a staging-signal heuristic, not a
cryptographic defeat of threshold secret-sharing (Section 8).

\textbf{Defense-in-depth across mechanism classes is required.} Even
with the strongest proxy at 38\%, and across all three tools, \textbf{15
of 24 vectors are covered by at least one tool; 9 are covered by none}
--- the registry/supply-chain surface (package squatting, provenance-gap
supply-chain poisoning, server impersonation), OS-isolation-dependent
vectors (sandbox escape), and semantic vectors (tool-transfer,
false-error escalation, configuration drift, command injection, consent
fatigue). These are not addressable by \emph{any} runtime proxy; they
require different mechanisms --- cryptographic attestation, OS
sandboxing, and client-side controls. A proxy is necessary but not
sufficient.

\textbf{Evasion robustness.} The corpus includes three obfuscated
variants of known attacks --- zero-width / bidirectional-control
characters, Cyrillic homoglyph substitution, and base64-wrapped
payloads. The tools resist \emph{different} evasions: mcp-bastion
catches the zero-width/bidi variant (via a hidden-character rule) but
misses the others; pipelock catches the homoglyph and base64 variants
(via a normalization pipeline) but misses zero-width; mcp-firewall
catches none. No single tool survives all three, so the defense-in-depth
conclusion holds at the evasion layer as well as the vector layer.

\textbf{Zero false positives.} No tool flagged any benign control,
across all 33 cases and every feature iteration.

\subsection{7. Discussion}\label{discussion}

\textbf{Benchmark-guided improvement.} The benchmark initially scored
mcp-bastion at 9\%, covering only the tool-definition layer. It then
guided five rounds of proxy-native hardening, each re-measured at zero
false positives: response scanning (9\%→18\%: response/retrieval
injection, result-borne leakage), argument/schema validation (18\%→23\%:
parameter smuggling, validation bypass), transport hardening (23\%→30\%:
plaintext-transport warning and a DNS-rebinding Origin check that
\emph{enforces} by rejecting cross-origin requests to a loopback
listener), sensitive-argument plus least-privilege scanning (30\%→34\%:
cross-tool exfiltration, over-broad scopes), and cross-tool correlation
(34\%→38\%: ShareLock split poisoning, plus MSTI caught by existing
pinning). The benchmark did not just measure the tool; it identified
concrete, layer-specific gaps and verified each fix --- and then defined
the proxy's ceiling, since the remaining 9 vectors are architecturally
outside a proxy's reach.

\textbf{A living benchmark tracks new attacks.} Vectors 23--24 were
published (June 2026) \emph{after} the v0.2 rubric was frozen; adding
them exercised the benchmark's intended lifecycle. Each was verified
against its primary source, encoded as a matched malicious/benign pair,
and re-run against every tool --- which immediately exposed that
ShareLock was covered by none, driving the cross-tool-correlation fix
above. A benchmark that only measured a fixed attack set would have
missed the gap entirely; the value is in re-measuring as the threat
surface moves.

\textbf{Tool-class taxonomy (what we score vs.~what we track).} Not
every MCP defense is a content-scanning proxy, and scoring all of them
on this rubric would be misleading. We separate two classes.
\emph{Content-scanning proxies} (mcp-bastion, mcp-firewall, pipelock)
inspect tool definitions, calls, or results at runtime; the rubric
measures exactly what they inspect, so we \textbf{score} them.
\emph{Access-control and isolation gateways} (e.g.~MCPX/Lunar, IBM
ContextForge, Docker's isolation gateway) enforce identity, policy, and
sandboxing rather than scanning content; the rubric's content-borne
vectors would score them near zero not because they are weak but because
they operate on a different axis. We \textbf{track} these in the
landscape (Section 2) without assigning a score, and note that a
class-appropriate enforcement-and-isolation rubric is future work.
Reporting a misleading number would be worse than reporting none.

\textbf{Corpus-encoding sensitivity (a fairness hazard).} We observed
that scores are sensitive to how an attack is \emph{encoded}. An early
corpus expressed exfiltration with sanitized placeholder text; against
it, mcp-firewall scored 5\%. Adding realistic encodings (an
AWS-key-shaped secret in a tool result; an exfiltration POST to the
cloud-metadata address) raised it to 13\% --- not by changing the tool,
but by probing capabilities the first corpus never exercised. We
therefore treat \emph{tool-neutral, realistic, multi-encoding} fixtures
as a fairness requirement, and we caution that \textbf{absolute totals
are corpus- dependent; the per-vector matrix (which vectors a tool
covers) is the more reliable read.}

\textbf{Adapter faithfulness.} The matched benign control also guards
against adapter over-reach. In wiring pipelock (an egress firewall), an
initial adapter synthesized an egress URL from an \emph{inbound} local
host, causing pipelock to ``block'' both the malicious and benign case
--- surfaced immediately as a false positive and corrected by
restricting URL synthesis to outbound endpoints.

\subsection{8. Limitations and Threats to
Validity}\label{limitations-and-threats-to-validity}

\begin{itemize}
\tightlist
\item
  \textbf{Taxonomy is a synthesis, not a standard.} The 24-vector set
  consolidates several proposed taxonomies; it is one reviewable
  synthesis. Framework mappings are indicative, single-author, and
  pending a second-reviewer pass; OWASP Agentic ASI05--ASI10 labels are
  unconfirmed.
\item
  \textbf{Small, seeded corpus.} Most vectors have 1--2 fixtures.
  Absolute coverage numbers should be read as lower bounds and are
  corpus-dependent (Section 7).
\item
  \textbf{The ShareLock check is a heuristic, not a proof.} Cross-tool
  correlation flags the \emph{staging signal} --- coordinated
  \texttt{share}/\texttt{checksum}-style metadata across a tool set ---
  that the published attack uses. It is not a cryptographic defeat of
  threshold secret-sharing; a sufficiently disguised split payload that
  avoids that signal can still evade it. We report it as
  \texttt{detect}, not \texttt{enforce}.
\item
  \textbf{Partial tool interfaces.} For pipelock we drive its offline
  URL/egress scanner and its MCP-response injection scanner, but not its
  poisoned-tool-description check or its full DLP/tool-call HTTP eval
  endpoint; its number therefore remains a lower bound on its total
  capability.
\item
  \textbf{Configuration dependence.} Results depend on each tool's
  configured policy; we commit the exact configs, but different policies
  would yield different numbers.
\item
  \textbf{Scope and tool class.} We score locally-deterministic
  content-scanning proxies; cloud-model defenders and
  access-control/isolation gateways are out of scope for scoring
  (Sections 4.4, 7) --- the latter would score misleadingly low on a
  content-borne rubric and are tracked, not scored.
\end{itemize}

\subsection{9. Future Work}\label{future-work}

Expand the corpus to multiple tool-neutral fixtures per vector and
deepen the evasion-robustness set; drive pipelock's remaining interfaces
(poisoned-description and DLP/tool-call eval endpoint) and add further
defenders (including a fully-local mode of cloud-gated gateways, if
exposed); publish a public leaderboard; complete the second-reviewer
pass on framework mappings and confirm ASI labels; and build an
\textbf{attestation-conformance} suite to test tool-integrity/provenance
schemes (ETDI, the Trustworthy Registry composition) as a coupled
bench-plus-attestation artifact.

\subsection{10. Conclusion}\label{conclusion}

We introduced a layer-aware, framework-mapped crosswalk of 24 MCP attack
vectors --- mapped to NIST AI RMF, the OWASP LLM and Agentic Top 10s,
STRIDE, and the NSA MCP Security guidance --- and an open,
vendor-neutral benchmark that measures how much of that surface a
defensive proxy actually covers, driving each tool through its real
code. The strongest of three open-source proxies reaches 38\% coverage,
and 9 of 24 vectors are undefended by any measured tool --- the
registry/supply-chain, OS-isolation, and semantic vectors that lie
outside a proxy's reach. The evidence argues for layered defense across
distinct mechanism classes, and the benchmark demonstrably guided one
proxy from 9\% to 38\% at zero false positives --- including a
cross-tool-correlation check added in response to a June-2026 attack
that no measured tool had covered. It provides a reusable, honest
instrument --- with matched benign controls, a runtime-integrity rule,
and a reproducibility rule --- for the community to measure MCP
defenders and track their progress.

\subsection{Availability}\label{availability}

\begin{itemize}
\item
  \textbf{Benchmark (mcp-defense-bench):} rubric, corpus, adapters,
  runner, leaderboard. Rubric/data: CC-BY-4.0; harness: Apache-2.0.
\item
  \textbf{Reference proxy (mcp-bastion):} Apache-2.0; published on npm.
\item
  Evaluated third-party tools are used under their own licenses
  (mcp-firewall: AGPL-3.0; pipelock: Apache-2.0 core) and are not
  redistributed here.
\item
  \textbf{Preprint of this paper:} DOI
  \href{https://doi.org/10.6084/m9.figshare.32978657}{10.6084/m9.figshare.32978657}
\item
  \textbf{mcp-defense-bench:}
  https://github.com/Gowthaman90/mcp-defense-bench · archived on Zenodo,
  DOI
  \href{https://doi.org/10.5281/zenodo.21346206}{10.5281/zenodo.21346206}
\item
  \textbf{mcp-bastion:} https://github.com/Gowthaman90/mcp-bastion ·
  npm: \texttt{mcp-bastion}
\end{itemize}

\subsection{References}\label{references}

\emph{Titles, authors, and identifiers below were verified against
primary sources on 2026-07-14. Each arXiv item resolves at
\texttt{https://arxiv.org/abs/\textless{}id\textgreater{}}. The
2603/2604/2606-series works are recent, non-peer-reviewed preprints and
may have updated versions; verify current version and page/DOI details
at camera-ready time.}

\begin{itemize}
\tightlist
\item
  {[}SoK-2512.08290{]} S. Gaire, S. Gyawali, S. Mishra, S. Niroula, D.
  Thakur, U. Yadav. ``Systematization of Knowledge: Security and Safety
  in the Model Context Protocol Ecosystem.'' arXiv:2512.08290,
  Dec.~2025.
\item
  {[}Formal-2604.05969{]} N. Acharya, G. K. Gupta. ``A Formal Security
  Framework for MCP-Based AI Agents: Threat Taxonomy, Verification
  Models, and Defense Mechanisms.'' arXiv:2604.05969, Apr.~2026.
\item
  {[}STRIDE-2603.22489{]} C. Huang, X. Huang, N. P. Tran, A. Milani
  Fard. ``Model Context Protocol Threat Modeling and Analyzing
  Vulnerabilities to Prompt Injection with Tool Poisoning.''
  arXiv:2603.22489, Mar.~2026.
\item
  {[}MCP-DPT-2604.07551{]} M. Rostamzadeh, S. Narula, N. Birhan, M.
  Ghasemigol, D. Takabi. ``MCP-DPT: A Defense-Placement Taxonomy and
  Coverage Analysis for Model Context Protocol Security.''
  arXiv:2604.07551, Apr.~2026.
\item
  {[}MSB-2510.15994{]} D. Zhang, Z. Li, X. Luo, X. Liu, P. Li, W. Xu.
  ``MCP Security Bench (MSB): Benchmarking Attacks Against Model Context
  Protocol in LLM Agents.'' arXiv:2510.15994, Oct.~2025 (ICLR 2026).
\item
  {[}MCPTox-2508.14925{]} Z. Wang, Y. Gao, Y. Wang, S. Liu, H. Sun, H.
  Cheng, G. Shi, H. Du, X. Li. ``MCPTox: A Benchmark for Tool Poisoning
  Attack on Real-World MCP Servers.'' arXiv:2508.14925, Aug.~2025.
\item
  {[}MCPBench-2508.20453{]} Z. Wang, Q. Chang, H. Patel, S. Biju, C.-E.
  Wu, Q. Liu, A. Ding, A. Rezazadeh, A. Shah, Y. Bao, E. Siow.
  ``MCP-Bench: Benchmarking Tool-Using LLM Agents with Complex
  Real-World Tasks via MCP Servers.'' arXiv:2508.20453, Aug.~2025.
\item
  {[}ETDI-2506.01333{]} M. Bhatt, V. S. Narajala, I. Habler. ``ETDI:
  Mitigating Tool Squatting and Rug Pull Attacks in Model Context
  Protocol (MCP) by using OAuth-Enhanced Tool Definitions and
  Policy-Based Access Control.'' arXiv:2506.01333, Jun.~2025.
\item
  {[}TrustReg-FI2026{]} L. Mas, J. Vilaplana, J. Rius, R. Spaimoc, J.
  Mateo. ``The Trustworthy Model Context Protocol (MCP) Registry: An
  Architectural Blueprint for Cryptographic Provenance and Runtime
  Integrity.'' Future Internet, vol.~18, no. 5, art. 243, 2026.
  doi:10.3390/fi18050243.
\item
  {[}WebMCP-MSTI-2606.06387{]} ``WebMCP: Tool-Surface Poisoning and
  Mid-Session Tool Injection in Browser-Hosted Model Context Protocol.''
  arXiv:2606.06387, Jun.~2026. (Vector 23.)
\item
  {[}ShareLock-2606.27027{]} ``ShareLock: A Stealthy Multi-Tool
  Threshold Poisoning Attack Against the Model Context Protocol.''
  arXiv:2606.27027, Jun.~2026. (Vector 24.)
\item
  {[}NIST-AI-RMF{]} National Institute of Standards and Technology (U.S.
  Dept. of Commerce). ``AI Risk Management Framework (AI RMF 1.0).''
  NIST AI 100-1, 2023 (functions: Govern, Map, Measure, Manage).
  https://www.nist.gov/itl/ai-risk-management-framework ·
  https://nvlpubs.nist.gov/nistpubs/ai/NIST.AI.100-1.pdf
\item
  {[}OWASP-LLM-2025{]} OWASP GenAI Security Project. ``OWASP Top 10 for
  LLM Applications,'' 2025 edition. https://genai.owasp.org/llm-top-10/
\item
  {[}OWASP-ASI-2026{]} OWASP GenAI Security Project. ``OWASP Top 10 for
  Agentic Applications,'' 2026 (published Dec.~9, 2025; categories
  ASI01--ASI10).
  https://genai.owasp.org/resource/owasp-top-10-for-agentic-applications-for-2026/
\item
  {[}NSA-MCP-2026{]} National Security Agency, AI Security Center. ``MCP
  Security: Security Design Considerations for AI-Driven Automation''
  (Cybersecurity Information Sheet, U/OO/6030316-26), May 2026.
  Recommendation areas: Authentication \& Access Control, Monitoring \&
  Logging, Data Classification \& Segmentation, Runtime Controls.
  https://www.nsa.gov/Portals/75/documents/Cybersecurity/CSI\_MCP\_SECURITY.pdf
\item
  {[}STRIDE{]} Microsoft. ``STRIDE --- Threat Modeling.''
  https://learn.microsoft.com/en-us/azure/security/develop/threat-modeling-tool-threats
\item
  {[}AgentDefense-Bench{]} A. Sanna. ``AgentDefense-Bench'' (open-source
  dataset), GitHub.
\item
  {[}Snyk-agent-scan{]} Snyk. ``agent-scan'' (formerly Invariant Labs
  MCP-Scan), GitHub / vendor tooling.
\end{itemize}

\end{document}
